% !TeX root = ../main.tex
\section{数式・表・グラフ}
\label{sec:results}
説明用に，抵抗 \(R=\SI{1.00}{\kilo\ohm}\) に電圧 \(V\) を加えた場合を考える。
オームの法則より電流 \(I\) は次のように求まる。
\begin{equation}
  I = \frac{V}{R}
  \label{eq:ohm}
\end{equation}
式\eqref{eq:ohm}による計算値を表\ref{tab:values}，
グラフを図\ref{fig:curve}に示す。
グラフは PGFPlots を用いて文書内で描く\cite{pgfplots}。

\begin{table}[htbp]
  \centering
  \caption{説明用の電圧と電流の計算値}
  \label{tab:values}
  \begin{tabular}{S[table-format=1.1] @{\hspace{12mm}} S[table-format=1.2]}
    \toprule
    {電圧 / \si{\volt}} & {電流 / \si{\milli\ampere}} \\
    \midrule
    0.0 & 0.00 \\
    1.0 & 1.00 \\
    2.0 & 2.00 \\
    3.0 & 3.00 \\
    4.0 & 4.00 \\
    5.0 & 5.00 \\
    \bottomrule
  \end{tabular}
\end{table}

\begin{figure}[htbp]
  \centering
  \begin{tikzpicture}
    \begin{axis}[
      width=0.78\linewidth,height=48mm,
      xlabel={Voltage / V},ylabel={Current / mA},
      xmin=0,xmax=5,ymin=0,ymax=5.5,
      grid=major,
      tick label style={font=\small},
      label style={font=\small}
    ]
      \addplot[color=accent,thick,mark=*]
        table[x=voltage,y=current,col sep=comma]{data/ohm.csv};
    \end{axis}
  \end{tikzpicture}
  \caption{抵抗値を一定としたときの電圧と電流の関係}
  \label{fig:curve}
\end{figure}
